\begin{abstract}
\bi{
LLM-based data agents automate complex analytics and predictive modeling over enterprise databases, yet independent sessions repeatedly explore the same tables to reconstruct identical domain computations.
Turning an agent's discovery into shared knowledge requires abstracting session-specific queries into generalizable definitions without dropping constraints, matching ambiguous natural-language requests without false hits, and keeping definitions valid when data updates violate relational assumptions under an unchanged schema.
We present \system, a shared layer between data agents and the DBMS that extracts, looks up, and maintains executable database knowledge.
From verified sessions, \system\ extracts parameterized metric and feature definitions together with structural and probe-derived data conditions, admitting a candidate only if its compiled SQL reproduces the verified answer.
For lookup, \system\ leverages each definition's explicit semantic choices to enumerate request intents and rule-label LLM-paraphrased training pairs for a lightweight cross-encoder judge; online, it serves confident matches directly and surfaces structural differences for homonymous or underspecified requests without per-lookup LLM calls.
As writes arrive, \system\ filters unaffected checks via table version stamps and delta rows, shares validation outcomes across definitions and agents within MVCC snapshots, and repairs broken definitions only when a unique single edit restores validity and reproduces the learning answer.
Evaluated on PostgreSQL across five analytical schemas and RelBench against eight baselines, \system\ reduces token usage and end-to-end latency by up to \PpMavraLateTokXNone$\times$ and \PpMavraLateSpeedupNone$\times$, answers 77\% of first-time lookups directly at 98.6\% precision, and reaches \CsBreakMavra\% accuracy with zero wrong answers under condition-breaking updates.
}
{基于大语言模型的数据智能体能够自动完成企业数据库上的复杂分析与预测建模，但相互独立的会话往往反复探查同一批数据表以重建相同的领域计算。要将单个智能体的发现转化为可共享的知识，需要将特定会话的查询无损抽象为可泛化的定义，在不产生误命中的前提下匹配含糊的自然语言请求，并在表结构不变而底层数据更新破坏关系假设时持续维护定义的有效性。本文提出位于数据智能体与 DBMS 之间的共享层 \system{}，用于提取、查找并维护可执行的数据库知识。\system{} 从核验会话中提取参数化的指标与特征定义，并推导结构与探查层面的数据条件，仅在编译生成的 SQL 复现核验答案时准予发布。在查找上，\system{} 利用定义中显式的语义口径枚举请求意图，对 LLM 一次性改写生成的训练样本按一致性规则自动打标，训练轻量级交叉编码器判别器；在线阶段无需调用 LLM，即可直接响应高置信匹配，并为同名或口径未明的请求返回结构化差异。在维护上，\system{} 按表版本戳与增量行两级过滤无关检查，在 MVCC 快照内跨定义和智能体共享验证结果，且仅当唯一单步修改恢复全部条件并复现学习答案时执行自动修复。在 PostgreSQL 上跨五种分析模式与 RelBench 同八种基线的对比表明，\system{} 将 token 消耗与端到端延迟最高降低 \PpMavraLateTokXNone{} 倍与 \PpMavraLateSpeedupNone{} 倍，以 98.6\% 的准确率直接回答 77\% 的首次查找，并在条件破坏更新下以零错答取得 \CsBreakMavra\% 的准确率。}
\end{abstract}
% The PVLDB template has no keyword block (ICDE keywords were: data agents, agent memory,
% semantic caching, knowledge reuse).
